% XeLaTeX 用プリアンブル（pandoc の header-includes として読み込む）
% 日本語専用パッケージ（xeCJK / luatexja）なしで組めるよう、
% 本文フォントに和欧両対応の Noto Serif CJK JP を使い、
% 和文の行分割は XeTeX 組み込みのロケール改行に任せる。
\XeTeXlinebreaklocale "ja"
\XeTeXlinebreakskip=0pt plus 0.5pt
\XeTeXlinebreakpenalty=0
\usepackage{amsmath,amssymb}
\setlength{\parindent}{1em}
\setlength{\parskip}{0.4\baselineskip}
\renewcommand{\arraystretch}{1.25}
\usepackage{etoolbox}
\AtBeginEnvironment{longtable}{\small}

% Noto Serif CJK JP に無い記号（上付き・下付きの添字文字、ℝ・ℕ、⟨ ⟩ など）用の補助フォント。
% 該当箇所は latex/fallback-glyphs.lua が {\fallbackfont ...} で包む。
\newfontfamily\fallbackfont{DejaVu Serif}
